\documentclass[11pt]{article}
\usepackage[margin=2.3cm]{geometry}
\usepackage{amsmath,amssymb}
\usepackage{xr}
\externaldocument{main}
\usepackage[colorlinks=true,linkcolor=blue,urlcolor=blue]{hyperref}
\usepackage{enumitem}
\usepackage{microtype}
\setlength{\parskip}{0.45em}
\setlength{\parindent}{0pt}
\newenvironment{comment}[1]{\par\medskip\noindent\textbf{Comment~#1.}\ \itshape}{\par}
\newcommand{\resp}{\par\noindent\textbf{Response.}\ }
\newcommand{\chg}{\par\noindent\textbf{Changes.}\ }

\title{Response to the reviewer (second round)\\[0.3em]\large A reliability-constrained, multi-objective
framework for minimum-mass design of steel frame structures under material and loading uncertainty}
\author{[Author name]}
\date{}

\begin{document}
\maketitle

We thank the reviewer for a second careful report. We agree that a response letter is not evidence
in itself. For each point below, we therefore name the new computation, the table or figure where
its result appears, and the sentences of the manuscript that changed. All new results were produced
by the seeded scripts in the repository, and the commit that produced them is given in the
manuscript (Data and code availability). Section, table and figure numbers refer to the revised
manuscript.

\section*{Summary of the main changes}
\begin{enumerate}[leftmargin=1.5em,itemsep=0.1em]
\item The three-storey design was re-optimised twice. The first design imposes the targets on 95\%
lower confidence bounds; the second adds a margin informed by the full-model discrepancy. Both
designs were validated against full second-order re-analysis (Section~\ref{sec:strict}).
\item The paper now distinguishes three claims everywhere: (P) the point estimate reaches the
target, (I) the interval contains it, and (L) the lower bound exceeds it. ``Within'' results are no
longer described as meeting the target (Section~\ref{sec:mc}).
\item A confidence-bounded formulation, RBDO-C95, was run for all seven frames. It is presented as
the recommended formulation when a minimum reliability is to be demonstrated
(Sections~\ref{sec:optimisation} and~\ref{sec:strict}, Table~\ref{tab:strict}).
\item The excluded part of the failure tail was tested with 3\,000 additional full re-analyses and
a residual-risk estimate. The statement that excluded samples ``cannot fail'' was removed
(Section~\ref{sec:fullmc}, Table~\ref{tab:tail}).
\item The full texts of Papadrakakis et al.\ (2004) and Truong and Kim (2017) were read, and
Table~\ref{tab:literature} and the novelty statement were revised accordingly. The full text of
Foley et al.\ (2007) could not be accessed.
\item Model-uncertainty-aware designs (RBDO-C-MU) quantify the extra steel needed to keep the target
when model factors are included (Table~\ref{tab:mudesign}).
\item An analysis of the effect of the 1\% drift bias of the $B_1$--$B_2$ surrogate on the drift
reliability was added (Table~\ref{tab:strict}).
\end{enumerate}

\section*{Major comments}

\begin{comment}{1 (three-storey design)}
The three-storey RBDO-C design ($\hat\beta=2.995$ [2.98, 3.01]) is accepted as ``within'', while the
full model gives $\Delta\beta=-0.07$ [$-0.17$, $-0.01$]. Re-optimise with a margin or a stricter
acceptance rule, report the extra mass, and revise the abstract and conclusions.
\end{comment}
\resp We did both, and validated the resulting designs in the full model. The results are in
Section~\ref{sec:strict} and Table~\ref{tab:validation}.
\begin{itemize}[itemsep=0.1em]
\item \emph{RBDO-C (8188~kg)}, the previous design. Full model: $\hat\beta^{full}_{ULS}=2.94$ [2.84,
3.01]. It is no longer described as meeting the target anywhere in the paper.
\item \emph{RBDO-C95 (8282~kg, $+1.15\%$)}, with the target imposed on the SAA lower confidence
bound. Surrogate: $3.068$ (lower bound 3.053), i.e.\ claim L. Full model: $3.01$ [2.89, 3.10] with
$\Delta\beta=-0.09$ [$-0.21$, $-0.01$]. The point estimate reaches the target, but the interval
contains it (claim P, I). The discrepancy is therefore systematic for this frame.
\item \emph{RBDO-C95 with a margin (8408~kg, $+2.69\%$).} The ULS target is raised to 3.21, i.e.\ by
the upper end of the discrepancy interval. Because the catalogue is discrete, the design jumps to
$\hat\beta_{ULS}=3.62$ [3.58, 3.66]. Full model: $3.61$ [3.45, 3.73], or 3.52 with the heavy-tailed
residual of the excluded samples. This meets the target with 95\% confidence on the sampled strata.
\end{itemize}
The steel needed to meet the ULS target with 95\% confidence in the full model is therefore 2.69\%
above RBDO-C and 5.17\% above DDO. The abstract, Discussion and Conclusions (items 1, 3 and 4) now
state this, and they describe the three-storey result as a limitation of the surrogate. We could not
isolate the cause of the discrepancy, and the manuscript says so rather than offering an
unverified explanation.

\begin{comment}{2 (acceptance rule)}
Distinguish the point estimate, the interval containing the target and the lower bound exceeding it;
present the confidence-bounded formulation as a substantive alternative; consider making it the
primary recommendation; use precise wording.
\end{comment}
\resp
\begin{itemize}[itemsep=0.1em]
\item \emph{Three claims.} Section~\ref{sec:mc} now defines the three claims (P), (I) and (L). It
states that only (L) establishes the target with 95\% confidence and that a ``within'' result is
inconclusive. The former sentence that accepted ``within'' as meeting the target was deleted.
\item \emph{Tables.} Table~\ref{tab:strict} classifies every design by the strongest claim, and
Table~\ref{tab:driftlimits} marks (L) in bold and (P/I) in italic.
\item \emph{Confidence-bounded formulation.} RBDO-C95 is defined in Section~\ref{sec:optimisation}
and was run for all seven frames. All its designs meet both targets with 95\% confidence in
independent verification. It costs 1.15--1.55\% more steel than RBDO-C in three frames and nothing in
three. In the 6-m bay frame it found a design 0.26\% lighter than the RBDO-C runs; we report this as
a limitation of the search budget. Against DDO, RBDO-C95 costs 1.1--11.7\%.
\item \emph{Recommendation.} The Discussion and Conclusion~4 now recommend RBDO-C95 whenever a
minimum reliability is to be demonstrated. RBDO-C is kept for the like-for-like comparison with
DDO, which is also judged on nominal values.
\item \emph{Wording.} The abstract and conclusions now state the claims separately: six of seven
RBDO-C designs reach (L), and all RBDO-C95 designs do.
\item \emph{Drift limits.} The comparison with tightened drift limits was redone for both claim levels
(Section~\ref{sec:targets}). Against deterministic designs with claim (L), RBDO-C95 is lighter in
three frames and 0.27\% (high wind) and 1.55\% (base) heavier in two.
\end{itemize}

\begin{comment}{3 (excluded tail)}
An observed maximum ratio of 1.049 is not a bound. Validate the threshold or estimate the residual
risk, explain the estimator, and check a lower threshold or a larger sample.
\end{comment}
\resp We agree, and the sentence that samples below 0.85 ``cannot fail'' has been removed. The
estimator is now described as resting on an assumption that is tested (Section~\ref{sec:fullmc}):
\begin{itemize}[itemsep=0.1em]
\item \emph{Lower threshold.} For every validated design, 300 samples of the stratum $0.70\le
U_s<0.85$ of the same pool were re-analysed with the full model, 3\,000 analyses in total. None
failed. The largest observed ratio is now described as an observed maximum: 1.060 in this stratum
and 1.061 over all 13\,042 re-analysed samples. Lowering the threshold to 0.70 leaves the point
estimates unchanged.
\item \emph{Limits of this check.} The 95\% intervals become much wider (lower ends 0.35--0.69
below the point estimates). This is because the stratum carries 6--37\% of the pool probability, and
300 failure-free samples cannot bound a probability of order $10^{-3}$ in it. We state this
explicitly; a distribution-free bound would need tens of thousands of full analyses per design.
\item \emph{Residual-risk estimate.} For all unanalysed samples with $U_s<0.85$, the residual is
estimated from a fitted model of $\ln(U_{full}/U_s)$: normal, and Student $t_3$ as a heavy-tailed
alternative. The normal model gives $<10^{-23}$. The $t_3$ model gives $2.7\times10^{-5}$ to
$1.8\times10^{-4}$, which lowers $\hat\beta^{full}$ by 0.01--0.02 for most designs, by 0.06 for nine
storeys and by 0.15 for 6-m bays. No target classification changes, except that the RBDO-C95
three-storey design drops from 3.01 to 2.97.
\end{itemize}
All of this is in Table~\ref{tab:tail}. The summary paragraph of Section~\ref{sec:fullmc} and
Conclusion~1 now say that the validation is statistical and that the tail beyond the observed
maximum is modelled, not observed.

\begin{comment}{4 (literature)}
Read the full texts of Papadrakakis et al.\ (2004), Foley et al.\ (2007) and Truong and Kim (2017);
check the four claimed contributions; update Table~1.
\end{comment}
\resp
\begin{itemize}[itemsep=0.1em]
\item \emph{Papadrakakis et al.\ (2004)}, full text read (the openly available PDF in the BazTech
repository). Two of our earlier implicit claims did not survive this reading, and we have corrected
them. They imposed the Eurocode 3/8 checks in both the deterministic and the reliability-based
problem, and they compared the two optima. The reliability-based design was heavier: their text
says about 26\%, while the volumes in their Table~6 differ by 34\%. Their probabilistic constraint
is on interstorey drift $>4\%$ of storey height under seismic loading, estimated with 5000 (MCS) or
1000 (LHS) samples, without confidence intervals. The introduction now states that imposing code
checks alongside a reliability constraint, and finding that the reliability-based design is
heavier, are \emph{not} new.
\item \emph{Truong and Kim (2017)}, full text read (an openly accessible copy of the article). Their
RBDO of steel space frames under gravity and wind also keeps deterministic strength, $H/400$ drift
and constructability constraints. Their failure criterion is load-carrying capacity from
plastic-hinge advanced analysis. Drift is not a probabilistic constraint. The sampling estimator is
verified against crude MCS on benchmark problems, but the failure probabilities of the final frame
designs are not re-estimated independently, and no confidence intervals are given.
\item \emph{Foley et al.\ (2007).} We could not access the full text: the publisher's page shows
only the abstract, and no open copy was found. We say so plainly. Its row in Table~1 and the rows of
the other studies we could assess only from abstracts are marked with $^\dagger$. The text makes
no novelty claim relative to them.
\item \emph{Revised contribution.} It is now stated relative to the two studies read in full. In
neither did we find: (i) a comparison under a code-calibrated gravity--wind model with deterministic
optima, including tightened drift limits and RBDO with and without code checks; (ii) a probabilistic
drift constraint next to the ULS constraint; (iii) confidence intervals or independent
re-estimation for the final designs; (iv) a tail validation of the surrogate against full
second-order re-analysis; or (v) a code release. Table~\ref{tab:literature} was updated with the
full-text information.
\end{itemize}

\begin{comment}{5 (model uncertainty)}
Make the conditional nature of the baseline explicit; avoid suggesting a real-world target is met;
explain the purpose of the model-factor analysis; if feasible, give a model-uncertainty-aware design.
\end{comment}
\resp
\begin{itemize}[itemsep=0.1em]
\item \emph{Conditional statements.} The conditional nature of the indices is now stated in
Section~\ref{sec:mc} (last paragraph), in Section~\ref{sec:sensitivity}, at the start of the
Conclusions and in the abstract. Section~\ref{sec:sensitivity} says explicitly that the baseline
designs do \emph{not} meet $\beta_T$ once model uncertainty with c.o.v.\ 0.10 is counted.
\item \emph{Purpose of the model factors.} Section~\ref{sec:probmodel} explains that they are used
both as a sensitivity study and, for two frames, inside the constraints.
\item \emph{Model-uncertainty-aware designs.} RBDO-C-MU (c.o.v.\ 0.05) was run for the base and
three-storey frames (Table~\ref{tab:mudesign}). Keeping $\beta_T=3.0$ with model uncertainty
included costs $+1.34\%$ (base, $\hat\beta_{ULS}=3.047$ with the factors) and $+2.69\%$ (three
storeys) relative to RBDO-C. Without the factors, the RBDO-C designs have only 2.77 and 2.70. Designs
for c.o.v.\ 0.10 were not computed; the manuscript says they would be heavier.
\end{itemize}

\begin{comment}{6 (cost model)}
Keep economy tied to mass or the illustrative index; do not imply validation against project
costs; state how the penalty ranges relate to plausible costs or that they are hypothetical.
\end{comment}
\resp Section~\ref{sec:costmodel} now states that $\tau$ and $\kappa$ are hypothetical sensitivity
parameters, not estimates of fabrication or connection costs, and that the ranking has not been
checked against project costs. The Discussion and Conclusion~5 limit all economic statements to steel
mass and the uncalibrated index.

\section*{Additional technical and editorial comments}

\begin{comment}{Table~3 vs.\ response letter (fabrication factor)}
\end{comment}
\resp Our first response letter was wrong on this point. The fabrication-factor model (normal, mean
1.00, c.o.v.\ 0.05) is taken from Ellingwood et al.\ (1980), as Table~\ref{tab:rv} states, and it
is sourced, not assumed. The assumed parameters are the mean-to-nominal ratio of $F_y$ (1.10), its
inter-member correlation (0.5), the out-of-plumbness standard deviation and the model factors. This
is what the ``Status'' column of Table~\ref{tab:rv} shows.

\begin{comment}{$B_1$--$B_2$ drift bias near the drift target}
\end{comment}
\resp A new computation increases all surrogate drifts by 1\% and re-evaluates $\beta_{drift}$ with
$10^6$ samples (Table~\ref{tab:strict}, last columns; stage \texttt{driftbias}). This lowers
$\beta_{drift}$ by 0.026--0.028. The base and high-wind RBDO-C designs (1.507, 1.511) would then fall
below 1.5, and the nine-storey design would become inconclusive. The RBDO-C95 designs keep claim L,
except the nine-storey design, which is identical to its RBDO-C design. Section~\ref{sec:strict}
states that drift targets met by less than about 0.03 are met only to the accuracy of the $B_1$--$B_2$
approximation.

\begin{comment}{Seeds and robustness}
\end{comment}
\resp Sections~\ref{sec:optimisation} and~\ref{sec:robust} and the ``demonstrated'' paragraph of the
Discussion describe the evidence as robustness within the tested search budget. The lighter 6-m bay
design found by the RBDO-C95 run is reported as an example of the limits of that budget.

\begin{comment}{Fully correlated loads and common imperfection}
\end{comment}
\resp A paragraph ``Dependence and spatial variability'' was added to Section~\ref{sec:probmodel}. It
states these assumptions and warns against extrapolating to buildings with substantial spatial
variability. The ``not demonstrated'' list in the Discussion repeats this.

\begin{comment}{Identifiable commit and executable workflow}
\end{comment}
\resp The Data and code availability statement now names the commit that produced the results of
this version. \texttt{reproduce.sh} and \texttt{jobs.txt} include every new run and stage
(confidence-bounded and model-uncertainty designs, \texttt{validateC}, \texttt{driftbias}). A
release tag and DOI will be created on acceptance.

\begin{comment}{Author placeholders}
\end{comment}
\resp The author names and affiliations remain placeholders in this version. They will be completed
before formal submission.

\begin{comment}{Headline claims}
\end{comment}
\resp The abstract and Conclusions were rewritten. They state the three claim levels separately,
give the three-storey full-model result, and state that all indices are conditional on the models.

\end{document}
